% Palett og TikZ-stiler for illustrasjonsslidene.
% Fargene er hentet fra den private pptx-decken (sam-semesteråpning).
% Hver tekstblokk er en egen node med font=-opsjon, slik at linjeavstanden
% følger fontstørrelsen (fontbytte inne i en node gir feil baselineskip).

\definecolor{illBg}{HTML}{F7F9FC}
\definecolor{illInk}{HTML}{10243E}
\definecolor{illGray}{HTML}{516274}
\definecolor{illPanel}{HTML}{E7EDF4}
\definecolor{illBlue}{HTML}{2F6FB0}
\definecolor{illTeal}{HTML}{2A9D8F}
\definecolor{illOrange}{HTML}{F4A261}
\definecolor{illRed}{HTML}{D65A5A}
\definecolor{illPurple}{HTML}{7566C8}

\newcounter{illnode}

% Kort med farget aksentstripe til venstre:
%   \illcard{(x,y)}{farge}{OVERLINJE}{Tittel}{brødtekst}
% Tegnes inne i en tikzpicture. Kortet er ca. 4.3 cm bredt.
\newcommand{\illcard}[5]{%
  \stepcounter{illnode}%
  \fill[#2] #1 rectangle ($#1+(0.09,-3.6)$);
  \node[anchor=north west, text width=3.9cm, inner sep=0pt,
        font=\scriptsize\bfseries, text=#2]
        (illn\theillnode a) at ($#1+(0.35,-0.05)$) {#3};
  \node[anchor=north west, text width=3.9cm, inner sep=0pt,
        font=\bfseries, text=illInk]
        (illn\theillnode b) at ($(illn\theillnode a.south west)+(0,-0.2)$) {#4};
  \node[anchor=north west, text width=3.9cm, inner sep=0pt,
        font=\scriptsize, text=illGray]
        at ($(illn\theillnode b.south west)+(0,-0.25)$) {#5};
}

% Mørk banner med hvit tekst over hele tekstbredden:
%   \illbanner{tekst}
\newcommand{\illbanner}[1]{%
  \begin{tikzpicture}
    \node[fill=illInk, rounded corners=5pt, text=white, font=\bfseries,
          minimum width=\textwidth, minimum height=0.95cm,
          text width=0.92\textwidth, align=center, inner ysep=6pt] {#1};
  \end{tikzpicture}}

% Trinnkort i prosesskjeden:
%   \illstep{(x,y)}{nr}{Tittel}{spørsmål}         (uten fyll)
%   \illstepx{(x,y)}{nr}{Tittel}{spørsmål}        (uthevet, lys blå plate)
\newcommand{\illstep}[4]{%
  \stepcounter{illnode}%
  \node[anchor=north west, text width=2.6cm, inner sep=0pt,
        font=\small\bfseries, text=illBlue]
        (illn\theillnode a) at ($#1+(0.12,-0.15)$) {#2};
  \node[anchor=north west, text width=2.6cm, inner sep=0pt,
        font=\bfseries, text=illInk]
        (illn\theillnode b) at ($(illn\theillnode a.south west)+(0,-0.18)$) {#3};
  \node[anchor=north west, text width=2.6cm, inner sep=0pt,
        font=\footnotesize, text=illGray]
        at ($(illn\theillnode b.south west)+(0,-0.22)$) {#4};
}
\newcommand{\illstepx}[4]{%
  \fill[illPanel, rounded corners=6pt] ($#1+(-0.18,0.35)$) rectangle ($#1+(2.95,-3.3)$);
  \illstep{#1}{#2}{#3}{#4}%
}
% Pil mellom trinnkort: \illarrow{(x,y)}
\newcommand{\illarrow}[1]{%
  \fill[illGray!70] ($#1+(0,0.09)$) -- ($#1+(0.34,-0.07)$) -- ($#1+(0,-0.23)$) -- cycle;
}

% E-rad for eksponeringskategoriene:
%   \illerow{(x,y)}{farge}{E0}{Skår: 0}{Tittel}{beskrivelse}
\newcommand{\illerow}[6]{%
  \fill[#2] #1 rectangle ($#1+(0.09,-1.55)$);
  \node[anchor=north west, inner sep=0pt, font=\large\bfseries, text=#2]
        at ($#1+(0.35,-0.02)$) {#3};
  \node[anchor=north west, inner sep=0pt, font=\footnotesize\bfseries, text=illGray]
        at ($#1+(1.35,-0.09)$) {#4};
  \node[anchor=north west, inner sep=0pt, font=\bfseries, text=illInk]
        at ($#1+(3.1,-0.05)$) {#5};
  \node[anchor=north west, text width=7.6cm, inner sep=0pt,
        font=\footnotesize, text=illGray]
        at ($#1+(0.35,-0.75)$) {#6};
}
